\documentclass[a4paper,11pt]{article}

% ASTRA Results template.
% Compatible target: Tectonic first, then latexmk/pdflatex fallback.

\usepackage[utf8]{inputenc}
\usepackage[T1]{fontenc}
\usepackage[french]{babel}
\usepackage{lmodern}
\usepackage{microtype}
\usepackage{amsmath}
\usepackage{amssymb}
\usepackage{booktabs}
\usepackage{longtable}
\usepackage{array}
\usepackage{geometry}
\usepackage{xcolor}
\usepackage{hyperref}

\geometry{margin=2.5cm}
\hypersetup{
  colorlinks=true,
  linkcolor=blue!45!black,
  urlcolor=blue!45!black,
  citecolor=blue!45!black,
  pdftitle={ASTRA Results},
  pdfauthor={ASTRA / Representations Procedurales Adressables}
}

\setlength{\parindent}{0pt}
\setlength{\parskip}{0.65em}
\renewcommand{\arraystretch}{1.15}

\newcommand{\AstraProject}{ASTRA / Repr\'esentations Proc\'edurales Adressables}
\newcommand{\AstraDecision}[1]{\textbf{\texttt{\detokenize{#1}}}}
\newcommand{\AstraMetric}[2]{\texttt{\detokenize{#1}} & #2\\}

\title{\textbf{ASTRA Results}\\Titre du rapport}
\author{\AstraProject}
\date{\today}

\begin{document}
\maketitle

\begin{abstract}
R\'esum\'e court du rapport Results. Le rapport Markdown d'analyse reste la
trace vivante; ce fichier \LaTeX{} est le livrable final fig\'e.
\end{abstract}

\section{Position}

D\'ecrire la position de l'\'etape dans ASTRA.

\section{R\'esultats}

\begin{longtable}{@{}p{0.52\linewidth}r@{}}
\toprule
M\'etrique & Valeur\\
\midrule
\AstraMetric{metric_name}{0}
\bottomrule
\end{longtable}

\section{D\'ecision}

\AstraDecision{RECALIBRATE}

\section{Limites}

Lister les limites sans sur-revendication scientifique.

\end{document}
